\documentclass[fleqn]{article}
\oddsidemargin 0.0in
\textwidth 6.0in
\thispagestyle{empty}
\usepackage{import}
\usepackage{amsmath}
\usepackage{graphicx}
\usepackage{flexisym}
\usepackage{calligra}
\usepackage{amssymb}
\usepackage{bigints} 
\usepackage[english]{babel}
\usepackage[utf8x]{inputenc}
\usepackage{float}
\usepackage[colorinlistoftodos]{todonotes}


\DeclareMathAlphabet{\mathcalligra}{T1}{calligra}{m}{n}
\DeclareFontShape{T1}{calligra}{m}{n}{<->s*[2.2]callig15}{}
\newcommand{\scriptr}{\mathcalligra{r}\,}
\newcommand{\boldscriptr}{\pmb{\mathcalligra{r}}\,}

\definecolor{hwColor}{HTML}{442020}

\begin{document}

  \begin{titlepage}

    \newcommand{\HRule}{\rule{\linewidth}{0.5mm}}

    \center

    \begin{center}
      \includegraphics[height=11cm, width=11cm]{asu.png}
    \end{center}

    \vline

    \textsc{\LARGE Statistical/Thermal Physics}\\[1.5cm]

    \HRule \\[0.5cm]
    { \huge \bfseries Homework 11}\\[0.4cm] 
    \HRule \\[1.0cm]

    \textbf{Behnam Amiri}

    \bigbreak

    \textbf{Prof: Michael Treacy}

    \bigbreak

    \textbf{{\large \today}\\[2cm]}

    \vfill

  \end{titlepage}

  \begin{enumerate}
    \item \textbf{5.81} Derive a formula, similar to equation $5.90$, for the shift in the freezing 
    temperature of a dilute solution. Assume that the solid phase is pure solvent, no solute. You should
    find that the shift is negative: The freezing temperature of a solution is \emph{less} than that
    of the pure solvent. Explain in general terms why the shift should be negative.

      % \textcolor{hwColor}{
      %   \\
      % }

    \item \textbf{5.82} Use the result of the previous problem to calculate the freezing temperature
    of seawater.

      % \textcolor{hwColor}{
      %   \\
      % }
    
    \item \textbf{6.2} Prove that the probability of finding an atom in any particular energy level is 
    $P(E)=\dfrac{1}{Z} e^{-F/kT}$, where $F=E-TS$ and the "entropy" of a level is $k$ times logarithm
    of the number of degenerate states for that level.

      % \textcolor{hwColor}{
      %   \\
      % }

    \item \textbf{6.3} Consider a hypothetical atom that has two states. a ground state 
    with energy zero and an excited state with energy $2 ~ eV$. Draw a graph...

      % \textcolor{hwColor}{
      %   \\
      % }

    \item \textbf{6.5} Imagine a particle that can be in only three states, with energies...
    \begin{enumerate}
      \item Calculate the partition function for this particle.

        % \textcolor{hwColor}{
        %   \\
        % }

      \item Calculate the probability for this particle to be in each of the three states.

        % \textcolor{hwColor}{
        %   \\
        % }

      \item Because the zero point for measuring energies is arbitrary, we...

        % \textcolor{hwColor}{
        %   \\
        % }

    \end{enumerate}

    \item \textbf{6.9} In the numerical example in the text, I calculated only the ratio of the 
    probabilities...
      \begin{enumerate}
        \item Estimate the partition function for a hydrogen atom at $5800 ~ K$...

          % \textcolor{hwColor}{
          %   \\
          % }

        \item Show that if all bound states are included in the sum, then the partition function...

          % \textcolor{hwColor}{
          %   \\
          % }

        \item When a hydrogen atom is in energy level $n$, the approximate radius of the electron...

          % \textcolor{hwColor}{
          %   \\
          % }

      \end{enumerate} 

    \item \textbf{6.23} For a $CO$ molecule, the constant $\epsilon$ is approximately $0.00024 ~ eV$...

      % \textcolor{hwColor}{
      %   \\
      % }

    \item \textbf{6.24} For an $O_2$ molecule, the constant $\epsilon$ is approximately $0.00018 ~ eV$. Estimate
    the rotational partition function for an $O_2$ molecule at room temperature.

      % \textcolor{hwColor}{
      %   \\
      % }

  \end{enumerate}

\end{document}
